\documentclass[a4paper,12pt]{article}
\usepackage[utf8]{inputenc}
\usepackage[T1]{fontenc}
\usepackage[french]{babel}
\usepackage{geometry}
\geometry{left=2.5cm, right=2.5cm, top=2.5cm, bottom=2.5cm}

\title{\textbf{Exercice : Sécurisation réseau TPE (Techno Vert)}}
\author{Maël - BTS SIO}
\date{\today}

\begin{document}

\maketitle

\section{Analyse du réseau actuel}

\subsection{Identification des risques et failles}
Voici 5 failles majeures identifiées dans l'infrastructure actuelle :

\begin{enumerate}
    \item \textbf{Absence de segmentation (Réseau plat) :} Tous les PC et le NAS sont sur le même réseau. Si un poste est infecté, le virus peut attaquer directement le NAS et les autres PC.
    \item \textbf{Usage dangereux du Wi-Fi :} L'utilisation du réseau "Invité" pour le travail interne expose les données de l'entreprise à n'importe quel visiteur connecté.
    \item \textbf{Matériel inadapté :} Le switch non-manageable empêche la création de VLANs pour séparer les flux.
    \item \textbf{Données exposées :} Le NAS contenant les données sensibles est accessible sans filtrage réseau spécifique.
    \item \textbf{Absence de visibilité :} Il n'y a aucune supervision pour détecter une attaque ou une panne en temps réel.
\end{enumerate}

\subsection{Limites de la Box Internet}
Une box grand public est insuffisante pour une entreprise de 20 personnes car :
\begin{itemize}
    \item Elle ne permet pas de créer des règles de pare-feu fines (filtrage par IP, port, protocole) entre différents zones locales.
    \item Elle ne gère pas (ou très mal) les VLANs nécessaires à la sécurité.
    \item Ses capacités de journalisation (logs) sont inexistantes, empêchant l'analyse post-incident.
\end{itemize}

\section{Conception d'une architecture sécurisée}

\subsection{Schéma réseau logique proposé}

\textbf{Équipements requis :}
\begin{itemize}
    \item 1 Pare-feu/Routeur professionnel (ex: Fortinet, Stormshield, pfSense).
    \item 1 Switch Manageable (Niveau 2/3).
    \item Points d'accès Wi-Fi professionnels (compatibles VLAN).
\end{itemize}

\textbf{Architecture des zones (VLANs) :}

\begin{itemize}
    \item \textbf{VLAN 10 - Administration :} Postes de direction et comptabilité (Données critiques).
    \item \textbf{VLAN 20 - Employés :} Postes de travail standards.
    \item \textbf{VLAN 30 - Serveurs/NAS :} Zone isolée, accessible uniquement via des ports spécifiques (ex: SMB, HTTPS) depuis les VLAN 10 et 20.
    \item \textbf{VLAN 90 - Invités (Wi-Fi) :} Accès Internet uniquement, isolation totale du réseau interne.
\end{itemize}

\subsection{Justification de la segmentation}
La segmentation réseau est essentielle pour \textbf{limiter la surface d'attaque}. En séparant les services :
\begin{itemize}
    \item On empêche la propagation latérale d'un virus (si un PC employé est infecté, le NAS et la Compta restent protégés par le pare-feu).
    \item On optimise les performances en réduisant le domaine de diffusion (broadcast).
\end{itemize}


\section{Mesures de sécurité à mettre en place}

Voici 6 mesures prioritaires pour sécuriser Techno Vert:

\begin{enumerate}
    \item \textbf{Mise en place de VLANs :} Pour séparer physiquement et logiquement les flux de données (Invités vs Production).
    \item \textbf{Filtrage strict (Pare-feu) :} Bloquer tout trafic entrant non sollicité et limiter le trafic sortant au strict nécessaire (bloquer les sites malveillants).
    \item \textbf{Sécurisation du Wi-Fi (WPA3) :} Utiliser le protocole WPA3-Enterprise si possible, ou à minima WPA2/3 avec un mot de passe complexe, sur un SSID séparé du réseau invité.
    \item \textbf{Sauvegardes hors ligne :} Configurer une sauvegarde du NAS sur un support déconnecté (disque USB ou Cloud immuable) pour contrer les ransomwares.
    \item \textbf{Mise à jour des firmwares :} Maintenir à jour le routeur, le switch et le NAS pour combler les failles de sécurité connues.
    \item \textbf{VPN pour l'accès distant :} Remplacer tout accès direct (port forwarding) par un VPN sécurisé pour les employés en télétravail.
\end{enumerate}

\section{Plan d'action}

\textbf{Priorité 1 : Immédiat (Urgence)}
\begin{itemize}
    \item Interdire l'usage du Wi-Fi "Invité" pour l'accès aux données internes.
    \item Changer les mots de passe par défaut de tous les équipements (Box, NAS).
    \item Vérifier que le NAS a une sauvegarde fonctionnelle.
\end{itemize}

\textbf{Priorité 2 : Court terme (Investissement)}
\begin{itemize}
    \item Acquérir et installer un Pare-feu et un Switch manageable.
    \item Configurer les VLANs et migrer les postes informatiques.
    \item Mettre en place les règles de filtrage inter-VLAN.
\end{itemize}

\textbf{Priorité 3 : Moyen terme (Amélioration)}
\begin{itemize}
    \item Mettre en place un outil de supervision (monitoring) pour alerter en cas de panne.
    \item Former les employés aux bonnes pratiques (ne pas se connecter au Wi-Fi invité).
\end{itemize}

\end{document}